% Solution: Mas-Colell, Whinston & Green (1995) Exercise 8.B.4. Input by \solitem.
% New in the 50-item list. Check: check_A1.py.
\begin{exsol}{MWG Exercise 8.B.4}{262}{Advanced}{4}
\asked Prove that, in iterated deletion of strictly dominated strategies, the set of surviving
strategies does not depend on the order in which dominated strategies are deleted.

\method Strict dominance and its iterated deletion (MWG \S8.B, Defs.~8.B.2--8.B.4 and the
iterated procedure, pp.~236--241; Lecture~1, slides~23 and~29; Tadelis \S4.2, pp.~63--65). Two
facts carry the proof: (i) dominance is \emph{inherited} when the opponents' strategy sets shrink;
(ii) in a finite game, a dominated strategy has a dominator that is itself undominated, and such a
dominator is never deleted. The game is finite, as in MWG \S8.B (finite $S_i$); a final remark
shows why finiteness is needed.

\solution
\textbf{Setting.} $\Gamma = [I, \{S_i\}, \{u_i\}]$ with every $S_i$ finite. For a product set
$T = \prod_i T_i \subseteq S$, say that $s_i \in T_i$ is \emph{strictly dominated in $T$ by}
$t_i \in S_i$ if $u_i(t_i, s_{-i}) > u_i(s_i, s_{-i})$ for every $s_{-i} \in T_{-i}$; it is
\emph{dominated in $T$} if this holds for some $t_i \in T_i$.

A \emph{deletion sequence} is a chain $S = T^0 \supseteq T^1 \supseteq \dots \supseteq T^K$ of
product sets in which each $T^{k+1}$ is obtained from $T^k$ by deleting a non-empty set of
strategies (of one or several players) each of which is dominated in $T^k$, and which stops only
when $T^K$ contains no dominated strategy. ``Order of deletion'' means the choice of which
dominated strategies to delete at each step; the lecture's procedure (Lecture~1, slide~29) is the
sequence that deletes all of them at once. Since $S$ is finite and each step deletes something,
every deletion sequence ends after finitely many steps. Claim: all deletion sequences end at the
same $T^K$.

\textbf{Lemma 1 (inheritance).} If $X \subseteq Y$ (that is, $X_j \subseteq Y_j$ for all $j$),
$s_i \in X_i$, and $s_i$ is strictly dominated by $t_i$ in $Y$, then $s_i$ is strictly dominated
by $t_i$ in $X$.

\emph{Proof.} The inequality $u_i(t_i, s_{-i}) > u_i(s_i, s_{-i})$ holds for every
$s_{-i} \in Y_{-i}$, hence for every $s_{-i} \in X_{-i} \subseteq Y_{-i}$. $\square$

\textbf{Lemma 2 (an undominated dominator).} Fix a product set $X$ and a player~$i$. If
$s_i \in X_i$ is strictly dominated in $X$ by some strategy in $S_i$, then there is a $t_i^* \in
S_i$ that strictly dominates $s_i$ in $X$ and is itself not strictly dominated in $X$ by any
strategy in $S_i$.

\emph{Proof.} Let $D = \{t \in S_i : t$ strictly dominates $s_i$ in $X\}$, non-empty and finite.
``$r$ strictly dominates $t$ in $X$'' is irreflexive and transitive (if
$u_i(r,\cdot) > u_i(t,\cdot)$ and $u_i(t,\cdot) > u_i(s_i,\cdot)$ on $X_{-i}$, then
$u_i(r,\cdot) > u_i(s_i,\cdot)$ there), so it is a strict partial order, and the finite set $D$ has
an element $t_i^*$ not dominated by any other element of $D$. If some $r \in S_i$ strictly
dominated $t_i^*$ in $X$, then by transitivity $r$ would dominate $s_i$, so $r \in D$, a
contradiction. $\square$

\textbf{Lemma 3 (the end point of any sequence is contained in every set of any other
sequence).} Let $(T^k)_{k=0}^K$ and $(R^l)_{l=0}^L$ be two deletion sequences, with end points
$T = T^K$ and $R = R^L$. Then $T \subseteq R^l$ for every $l$; in particular $T \subseteq R$.

\emph{Proof}, by induction on $l$. For $l = 0$, $T \subseteq S = R^0$. Suppose $T \subseteq R^l$,
and suppose, for a contradiction, that some $s_i \in T_i$ is deleted at step $l+1$ of the second
sequence. Then $s_i$ is strictly dominated in $R^l$ by some $t_i \in R^l_i$. Since
$T \subseteq R^l$, Lemma~1 shows that $s_i$ is strictly dominated in $T$ by $t_i$. By Lemma~2
(with $X = T$) there is a $t_i^* \in S_i$ that strictly dominates $s_i$ in $T$ and is not strictly
dominated in $T$ by any strategy in $S_i$.

Now $t_i^* \in T_i$. Indeed, if $t_i^*$ had been deleted at some step $k+1$ of the first
sequence, it would be strictly dominated in $T^k$ by some $r \in T^k_i$; since $T \subseteq T^k$,
Lemma~1 makes $r$ strictly dominate $t_i^*$ in $T$, contradicting the choice of $t_i^*$.

So $s_i \in T_i$ is strictly dominated in $T$ by $t_i^* \in T_i$: $s_i$ is dominated in $T$. But
$T$ is the end point of a deletion sequence and contains no dominated strategy. This contradiction
shows that no strategy of $T$ is deleted at step $l+1$, so $T \subseteq R^{l+1}$. $\square$

\textbf{The theorem.} Apply Lemma~3 twice: $T \subseteq R$, and, exchanging the roles of the two
sequences, $R \subseteq T$. Hence $T = R$: every order of deletion ends at the same set, which is
the set $S^\infty$ produced by deleting everything dominated in each round. $\blacksquare$

(No $T^k_i$ ever becomes empty: dominance in $T^k$ is a strict partial order on the finite set
$T^k_i$, so $T^k_i$ has an element that no other element of $T^k_i$ dominates, and that element is
not deleted at step $k+1$.)

\textbf{Why finiteness matters.} Take one player with $S = [0,1]$ and $u(s) = s$ for $s < 1$,
$u(1) = 0$. Every $s < 1$ is strictly dominated by $(1+s)/2$, and $s = 1$ by $s = \tfrac12$.
Deleting everything dominated at once leaves the empty set. Deleting only the interval $(0,1)$
first leaves $\{0, 1\}$, where $u(0) = u(1) = 0$ and nothing is dominated, so the process stops at
$\{0,1\}$. Lemma~2 is what fails: the set of dominators of a strategy has no undominated element.

\textbf{Why ``strictly'' matters} (the manual's note, made concrete).
\begin{center}
\begin{tabular}{r|c|c|}
\multicolumn{1}{r}{1 $\backslash$ 2} & \multicolumn{1}{c}{$L$} & \multicolumn{1}{c}{$R$}\\ \cline{2-3}
$T$ & $1,\ 1$ & $0,\ 0$ \\ \cline{2-3}
$M$ & $1,\ 1$ & $2,\ 1$ \\ \cline{2-3}
$B$ & $0,\ 0$ & $2,\ 1$ \\ \cline{2-3}
\end{tabular}
\end{center}
With weak dominance: deleting $T$ (weakly dominated by $M$) first makes $L$ weakly dominated by $R$
on rows $\{M,B\}$, and the process ends at $\{M,B\}\times\{R\}$ with payoffs $(2,1)$. Deleting $B$
(weakly dominated by $M$) first makes $R$ weakly dominated by $L$ on rows $\{T,M\}$, and it ends at
$\{T,M\}\times\{L\}$ with payoffs $(1,1)$. Lemma~1 fails for weak dominance: removing opponent
strategies can remove the only profile where the inequality was strict.

\interp IESDS is a well-defined solution concept: what common knowledge of rationality rules out
does not depend on the order in which an analyst (or the players) notices dominated strategies.
That is why the slides can delete everything at once and why textbook solutions that delete one
strategy at a time reach the same answer. Weak dominance does not have this property, which is one
reason the course does not iterate it.

\verify \texttt{check\_A1.py} draws 2{,}000 random finite games (two or three players, up to four
strategies each, small integer payoffs so that ties are common), runs 20 random deletion orders in
each (at each step a random non-empty subset of the currently dominated strategies is deleted) and
checks that all of them end at the maximal-deletion set; it also checks the weak-dominance example
(two orders, two different end points) and the one-player continuum example on a grid of exact
rationals. \emph{Comparison with the MWG solutions manual} (pp.~8-2--8-3): it gives the same key idea in
two lines (a dominated $s_i$ has a dominator that is not itself dominated; if $s_i$ is not deleted
now it is still dominated, and deleted, later) and the same remark that order matters for weak
dominance. It does not show that the undominated dominator survives every later round, nor where
finiteness enters; Lemmas~2--3 supply both. The conclusion agrees.

\mistakes
\begin{itemize}
  \item Arguing ``a dominated strategy stays dominated, so it is deleted eventually'' without
  checking that its dominator is still there. The dominator may be deleted; Lemma~2 picks one that
  cannot be.
  \item Proving only that the maximal procedure deletes more at each round. The claim is about
  the end points of \emph{every} order, including orders that delete one strategy at a time.
  \item Extending the result to weak dominance, or to infinite games without further
  assumptions; both counterexamples above show it fails.
\end{itemize}

\answer{In a finite game, strict dominance is inherited by smaller games (Lemma~1), and every
dominated strategy has a dominator that is itself undominated (Lemma~2), which no deletion
sequence ever removes. So the end point $T$ of any deletion sequence is contained in every stage
of any other sequence (Lemma~3), hence in its end point $R$; by symmetry $R \subseteq T$. All
orders give the same surviving set $S^\infty$. (False for weak dominance and, in general, for
infinite games: examples above.)}
\end{exsol}
